\section{Conclusion}

The United States interstate highway system is 48,800 miles of concrete and asphalt with no official tier classification. Every corridor competes as an equal for maintenance funding, capacity investment, and operational priority. The result is an investment process driven by political visibility rather than network function.

We have derived, from publicly available federal data, a four-tier arterial classification of the 227 US interstate corridors. Eight Primary Arteries --- I-80, I-90, I-10, I-40, I-5, I-95, I-35, I-75 --- carry the majority of national truck freight ton-miles and occupy the highest betweenness centrality positions in the national network. These are the load-bearing spines. Disruption on a Tier 1 corridor cascades nationally; disruption on a Tier 4 route is locally significant and nationally invisible.

The tier classification has three direct uses for Interstate 2.0 planning:

\textbf{Investment sequencing.} Tier 1 arterials should receive the full Interstate 2.0 feature set: managed freight lanes, intermodal integration, EV charging corridors, resilience hardening. Tier 2 connectors receive a subset. Tier 3 and 4 routes receive maintenance and basic access improvements. This is how transit agencies sequence capital investment; it should be how highway investment is sequenced too.

\textbf{Redundancy prioritization.} The most dangerous gaps in the national network are not missing Tier 3 routes but missing redundancy for Tier 1 routes. Donner Pass on I-80 has no Tier 1 parallel. I-10's Gulf Coast segments have no inland alternative. The tier classification surfaces these single points of failure in a way that route-by-route analysis cannot.

\textbf{Performance measurement.} Tier 1 arterials should be held to higher performance standards --- tighter PTI targets, higher reliability commitments, more frequent condition inspection. The current uniform standard obscures the asymmetric consequences of degradation on different tier corridors.

The schematic arterial map is the visible output of this work. Like the Beck diagram for the London Underground, it reduces a complex geographic network to a legible hierarchy. It does not replace detailed corridor analysis; it provides the strategic frame within which that analysis occurs. The country chose interstates. The least it can do is understand which ones matter most.

\subsection*{Future Work}

This paper scores corridors on static 2018 HPMS data. Dynamic tier classification --- updated annually as AADT and freight flows shift --- would provide a real-time picture of arterial stress. Integration with FAF5 routing estimates would allow tier assignments based on actual freight flows rather than proxy metrics. The tier classification also provides a natural input to the investment allocation LP described in \citet{route_invest_2026}: Tier 1 corridors receive constrained priority before Tier 2 corridors are funded.
